% Section 2.9, from lectures/L02/README.md: the objectives, the questions to test yourself, and the
% next lecture.

\needspace{10\baselineskip}
\section{Review}
\label{c2:sec:review}

\subsection*{What you should be able to do now}
After this chapter, you should be able to:
\begin{itemize}
  \item Derive a minimized equation from a Karnaugh map and turn it into a VHDL architecture.
  \item Read and write \code{std_logic_vector} ports and signals.
  \item Write a \code{process} with a correct sensitivity list and a \code{case} covering every input
    value.
  \item Build a design out of more than one entity with a positional \code{port map}, and say why
    instantiating a module twice puts two copies of its gates on the FPGA.
  \item Verify a module by running its testbench, and read a failed assertion.
\end{itemize}

\needspace{6\baselineskip}
\subsection*{Questions to test yourself}
You should be able to explain:
\begin{enumerate}
  \item Given a truth table, can you derive its equation both directly and via a Karnaugh map? What
    is different about the two results?
  \item Why can a \code{case} statement not appear directly in an architecture body?
  \item What does \code{when others} protect you against, given that a selector ``can only'' be
    \code{0} or \code{1}?
  \item In \code{display1: entity work.display}, what does \code{display1} name, and what does
    \code{work.display} name? Could a second instance reuse the same label?
  \item \code{hex_display}'s architecture contains no logic at all. Where do the gates in the finished
    design come from?
  \item You instantiated \code{display} twice rather than writing its \code{case} twice. Does the FPGA
    end up with one copy of that logic or two? So what did you actually gain?
  \item This book always writes \code{port map} positionally. What breaks if an entity's ports are
    declared in a different order from the one its testbench expects, and when would you find out?
\end{enumerate}

\needspace{6\baselineskip}
\subsection*{Looking ahead}
Chapter~\ref{c3:ch} is about sequential logic, circuits that remember:
\begin{itemize}
  \item The D latch, and the D flip-flop built from a pair of them.
  \item Clocking, and when a signal assignment actually takes effect.
  \item Edge detection: a button press that toggles an LED.
\end{itemize}
